\section{Quantum gates}
\label{sec:quantum-gates}
\label{sec:hamiltonians-to-gates}

Initialization and readout provide a known input and a way to
measure the final spin populations. Between these stages, the
Hamiltonians in Section~\ref{sec:spin-control-exchange} control the
coherent evolution of the qubits. Implementing a particular gate
requires choosing the pulse amplitudes, durations, phases, and
ordering so that the resulting evolution has the desired action
on the computational states. For a pulse described by $H(t)$ over
a time interval $\tau$, the evolution is $U(\tau)$ in
Eq.~\eqref{eq:gate-time-evolution}.

To implement a unitary gate, the pulse must return every initial
computational state to the chosen computational space. If
$P_{\mathrm q}$ projects onto that space, this condition is
$(I-P_{\mathrm q})U(\tau)P_{\mathrm q}=0$; the restriction of
$P_{\mathrm q}U(\tau)P_{\mathrm q}$ then gives the gate operator.
Intermediate admixture of charge or other spin states is allowed,
provided that no population remains outside the computational space
at the end of the pulse. A calculation based on an effective spin
Hamiltonian must account for any such intermediate excursions
through the approximations used to obtain it.

\subsection{Gates for single-spin qubits}
\label{sec:single-qubit-gates}

\paragraph{Transverse rotation gates: \texorpdfstring{$R_x$ and $R_y$}{Rx and Ry}.}
We want to choose the axis and angle of a logical spin rotation.
The drive phase sets the transverse axis, while its amplitude and
duration set the angle.
In the logical basis $(\ket0,\ket1)=(\ket{\downarrow},\ket{\uparrow})$,
the Pauli operators obey
$\sigma_x=X$, $\sigma_y=-Y$, and $\sigma_z=-Z$.

Writing Eq.~\eqref{eq:esr-propagator} in this basis and allowing
a time-dependent drive envelope $\Omega(t)$ gives
\begin{equation}
  H_{\mathrm{rot,q}}(t)=\frac{\hbar}{2}
       \left[-\delta Z+\Omega(t)(\cos\phi\,X-\sin\phi\,Y)\right].
  \label{eq:logical-phase-drive}
\end{equation}
To reach resonance, tune the drive frequency to the measured spin
resonance $\omega_Z$ (Appendix~\ref{sec:measurement-spin-spectroscopy}).
On resonance ($\delta=0$), fixing $\phi$ keeps the rotation axis constant even
if the envelope varies. The angle depends on the full pulse area,
including its rise and fall,
$\vartheta=\int_0^\tau\Omega(t)\,dt$:
\begin{equation}
  U_\phi(\vartheta)=
  \exp\!\left[-\frac{i\vartheta}{2}
                    (\cos\phi\,X-\sin\phi\,Y)\right].
  \label{eq:phase-controlled-rotation}
\end{equation}
For a constant-amplitude pulse, the rotation rate $\Omega$ at that
amplitude is measured (Appendix~\ref{sec:measurement-rabi}) and the
duration is chosen so that $\Omega\tau=\vartheta$. For a pulse with
finite rising and falling edges, choose the amplitude and hold time
from the calibrated area of the complete waveform
(Appendix~\ref{sec:calibration-spin-rotations}).

To select the transverse axis, calibrate the relative phase of
successive pulses against a common reference
(Appendix~\ref{sec:calibration-spin-corrections}).
Taking $\phi=0$ as the $x$ reference, the two phase settings give
\begin{equation}
 \begin{aligned}
  \delta=0,\quad\phi=0
    &: \qquad U_0(\vartheta)=R_x(\vartheta),\\
  \delta=0,\quad\phi=-\pi/2
    &: \qquad U_{-\pi/2}(\vartheta)=R_y(\vartheta),
 \end{aligned}
 \qquad \vartheta=\int_0^\tau\Omega(t)\,dt.
 \label{eq:single-spin-transverse-gates}
\end{equation}
\paragraph{Phase rotation gate: \texorpdfstring{$R_z$}{Rz}.}
A logical $z$ rotation changes the relative phase without
moving population between the basis states. At the chosen holding
voltages, measure the detuning from the control reference
(Appendix~\ref{sec:measurement-ramsey-echo}). With the transverse
drive off, choose the waiting time to accumulate the required angle:
\begin{equation}
  U_z(\tau)=\exp\!\left[\frac{iZ}{2}\int_0^\tau\delta(t)\,dt\right]
            =R_z\!\left(-\int_0^\tau\delta(t)\,dt\right).
  \label{eq:detuning-z-gate}
\end{equation}
For a constant measured detuning, the condition is
$\vartheta_z=-\delta\tau$. Phase-coherent control can also track
a $z$ rotation by updating the programmed phases of later pulses;
those phases use the same calibrated reference as the transverse
gates above.

The calibrated rotations implement the Pauli and Hadamard gates
through the sequences in Eq.~\eqref{eq:elementary-single-qubit-gates}.

\paragraph{Exchange gates: SWAP and \texorpdfstring{$\sqrt{\mathrm{SWAP}}$}{square-root SWAP}.}
\label{sec:exchange-two-qubit-gates}
We next use exchange to entangle two single-spin qubits and combine
it with local rotations to obtain CNOT.
Keep the dots in $(1,1)$ and remove common Zeeman evolution.
The local Zeeman-energy difference must be calibrated for the
intended exchange operating point
(Appendix~\ref{sec:calibration-two-spin-exchange}). Choose holding
and exchange voltages for which this difference has negligible
action during the pulse. If no such setting is available, the
measured difference must be included in the pulse Hamiltonian
instead of using the exchange-only form below.
The exchange pulse gives
\begin{equation}
  \begin{aligned}
  H_{\mathrm{ex}}(t)&=\frac{J(t)}{4}
                  (X_1X_2+Y_1Y_2+Z_1Z_2),\\
  U_{\mathrm{ex}}(\theta)&=
       e^{-i\theta(X_1X_2+Y_1Y_2+Z_1Z_2)/4},
  \qquad \theta=\frac{1}{\hbar}\int_0^\tau J(t)\,dt.
  \end{aligned}
  \label{eq:logical-exchange-gate}
\end{equation}
For the SWAP operator in Eq.~\eqref{eq:exchange-gate-matrices},
$\mathrm{SWAP}=(I+X_1X_2+Y_1Y_2+Z_1Z_2)/2$
and $\mathrm{SWAP}^2=I$,
so
\begin{equation}
  U_{\mathrm{ex}}(\theta)=e^{i\theta/4}
       \left[\cos(\theta/2)I-i\sin(\theta/2)\mathrm{SWAP}\right].
  \label{eq:partial-swap}
\end{equation}
To set $\theta$, calibrate the exchange strength and the phase
accumulated on the pulse edges along the compensated barrier-voltage
path (Appendix~\ref{sec:calibration-two-spin-exchange}). The barrier
pulse amplitude and hold time are then chosen to reach the target
area (Appendix~\ref{sec:calibration-partial-swap}).
The pulse areas $\theta=\pi$ and $\theta=\pi/2$ implement the
two gates in Eq.~\eqref{eq:exchange-gate-matrices}, up to global phase:
\begin{equation}
 \mathrm{SWAP}=e^{i\pi/4}U_{\mathrm{ex}}(\pi),\qquad
 \sqrt{\mathrm{SWAP}}=e^{i\pi/8}U_{\mathrm{ex}}(\pi/2).
 \label{eq:exchange-gate-pulse-areas}
\end{equation}
SWAP exchanges the qubits without entangling a product state.
The half-area pulse can entangle them; for example,
\begin{equation}
  U_{\mathrm{ex}}(\pi/2)\ket{01}
    =\frac{e^{i\pi/8}}{\sqrt{2}}(\ket{01}-i\ket{10}).
  \label{eq:exchange-entanglement}
\end{equation}
The displayed output is maximally entangled~\cite{Loss1998}.

\paragraph{Controlled gates: CZ and CNOT.}
To construct CZ from partial exchange, we need to cancel the
$X_1X_2$ and $Y_1Y_2$ contributions while retaining $Z_1Z_2$.
Set $G=U_{\mathrm{ex}}(\pi/2)$.
Conjugation by the local phase flip $Z_1$ reverses the first two
terms but leaves the third unchanged. Because the three terms
commute, two exchange pulses give
\begin{equation}
  U_{ZZ}\equiv G Z_1 G Z_1
       =\exp(-i\pi Z_1Z_2/4).
  \label{eq:exchange-zz-construction}
\end{equation}
The phase flip $Z_1$ is supplied by a physical local $\pi$ rotation
[Eq.~\eqref{eq:elementary-single-qubit-gates}], with exchange
suppressed during that rotation. Each $G$ pulse uses the calibrated
exchange area $\pi/2$.
The resulting $U_{ZZ}$ assigns phases to all four basis states.
Local $z$ rotations remove the unwanted single-qubit phases, leaving
a minus sign only on $\ket{11}$:
\begin{equation}
  \mathrm{CZ}
     =e^{-i\pi/4}R_{z,1}(-\pi/2)R_{z,2}(-\pi/2)U_{ZZ}.
  \label{eq:exchange-cz-construction}
\end{equation}
For the complete sequence, the conditional phase and the two local
phases must be measured
(Appendix~\ref{sec:calibration-single-spin-exchange-cz}). Adjust the
exchange hold times until the conditional phase is $\pi$ modulo
$2\pi$, and set the local $z$ corrections from the measured local
phases.
Applying the calibrated Hadamard gate to the target before and
after this CZ gives CNOT [Eq.~\eqref{eq:exchange-cnot-construction}].
This sequence requires separately addressable single-spin pulses,
with charge leakage and unwanted Zeeman terms suppressed during
exchange. It provides the entangling gate that, together with
arbitrary single-qubit rotations, forms a universal gate
set~\cite{Loss1998}.

\subsection{Gates for singlet--triplet qubits}
\label{sec:st-gates}

For the basis $(\ket S,\ket{T_0})$, we again need controls that
change both relative phase and population.
Exchange changes the singlet--triplet energy splitting and hence
their relative phase. A difference between the longitudinal fields
on the two dots couples the two states.
For local Zeeman energies $E_{Z,i}=g_i\mu_B B_{z,i}$,
with effective factors $g_i$ and fields $B_{z,i}$, the physical
Hamiltonian is
\begin{equation}
  H=J\,\mathbf{s}_1\cdot\mathbf{s}_2
    +E_{Z,1}s_{1z}+E_{Z,2}s_{2z}.
  \label{eq:st-physical-hamiltonian}
\end{equation}
This Hamiltonian preserves the $m=0$ computational space.
Exchange gives energies $-3J/4$ and $J/4$ to $\ket S$ and $\ket{T_0}$,
while $\Delta E_Z=E_{Z,1}-E_{Z,2}$ couples them.
With $Z_{\mathrm{ST}}=\ket S\bra S-\ket{T_0}\bra{T_0}$ and
$X_{\mathrm{ST}}=\ket S\bra{T_0}+\ket{T_0}\bra S$,
the logical Hamiltonian, after dropping $-J/4$, is
\begin{equation}
  H_{\mathrm{ST}}=-\frac{J}{2}Z_{\mathrm{ST}}
                   +\frac{\Delta E_Z}{2}X_{\mathrm{ST}},
  \label{eq:st-effective-hamiltonian}
\end{equation}
A uniform longitudinal field vanishes in the $m=0$ space; the
difference $\Delta E_Z$, rather than the common field, sets the
mixing term. Fluctuations of either $J$ or $\Delta E_Z$ therefore
change the logical evolution.

\paragraph{Logical \texorpdfstring{$R_x$ and $R_z$}{Rx and Rz} gates.}
For constant controls held for a time $\tau$, the two pure-axis
rotations are
\begin{equation}
  \begin{aligned}
   J=0 &: \quad U_{\mathrm{ST}}=R_{x,\mathrm{ST}}(\Delta E_Z\tau/\hbar),\\
   \Delta E_Z=0 &: \quad U_{\mathrm{ST}}=R_{z,\mathrm{ST}}(-J\tau/\hbar).
  \end{aligned}
  \label{eq:st-axis-limits}
\end{equation}
To obtain $R_{x,\mathrm{ST}}$, set the compensated barrier and detuning
voltages so that exchange is negligible, measure $\Delta E_Z$
there, and choose $\tau$ for the required angle
(Appendix~\ref{sec:calibration-st-rotations}). A direct
$R_{z,\mathrm{ST}}$ requires a setting with negligible
$\Delta E_Z$; at that setting, the measured exchange strength
determines the pulse duration for the required angle.

\paragraph{Rotations with a fixed field gradient.}
With a fixed nonzero $\Delta E_Z$, turning on $J$ tilts the axis
rather than giving a pure $z$ rotation. The values of $J$ and
$\Delta E_Z$ must be measured at each selected voltage setting
(Appendix~\ref{sec:calibration-st-rotations}). Holding both controls
constant for time $\tau$ gives the gate below, with
$E_{\mathrm{ST}}=\sqrt{J^2+(\Delta E_Z)^2}$:
\begin{equation}
  U_{\mathrm{ST}}(\tau)
   =\cos\!\left(\frac{E_{\mathrm{ST}}\tau}{2\hbar}\right)I
    -i\sin\!\left(\frac{E_{\mathrm{ST}}\tau}{2\hbar}\right)
       \frac{\Delta E_Z X_{\mathrm{ST}}-JZ_{\mathrm{ST}}}
            {E_{\mathrm{ST}}}.
  \label{eq:st-pulse-rotation}
\end{equation}
The rotation axis is proportional to $(\Delta E_Z,0,-J)$.
Choose voltage settings whose measured $J$ and $\Delta E_Z$ give
nonparallel axes. The durations of these voltage
holds can then be chosen to compose the desired logical rotation.
Such control has been demonstrated in a silicon double dot with a
micromagnet~\cite{Wu2014}.
If $J/\Delta E_Z$ changes within a pulse, the axis also changes
and time ordering is generally required.

\paragraph{Encoded CZ and CNOT gates.}
To entangle two singlet--triplet qubits, couple two double dots
$A,B$ containing four electrons in total.
Capacitive coupling makes the phase acquired by one qubit depend
on the state of the other, adding an interaction to the separate
single-qubit Hamiltonians:
\begin{equation}
 H_{\mathrm{logical}}
   =H_A\otimes I_B+I_A\otimes H_B+H_{\mathrm{cap}}.
 \label{eq:st-coupled-logical-hamiltonian}
\end{equation}
Here $H_A,H_B$ include local energy shifts and $H_{\mathrm{cap}}$
is the additional two-qubit interaction.
After eliminating charge excitations, assume this coupling is
diagonal in the product singlet--triplet basis. Here $\ket S$
denotes the charge-hybridized singlet. Let
$E_{00},E_{01},E_{10},E_{11}$ be the energies of $\ket{SS}$,
$\ket{ST_0}$, $\ket{T_0S}$, and $\ket{T_0T_0}$, respectively.
Its conditional part is fixed by
$\kappa=E_{00}-E_{01}-E_{10}+E_{11}$: independent shifts of
either qubit cancel in this combination. Absorbing those shifts
into $H_A,H_B$, and writing $P_{S,q}=(I+Z_q)/2$ with $Z_q$ the
encoded Pauli operator on $q=A,B$, gives the interaction
(Appendix~\ref{sec:calibration-capacitive-st-interaction})
\begin{equation}
 H_{\mathrm{cap}}=\kappa P_{S,A}P_{S,B}
     =\frac{\kappa}{4}(I+Z_A+Z_B+Z_AZ_B).
 \label{eq:st-capacitive-coupling}
\end{equation}
The Zeeman differences at the joint interaction setting must also
be measured. Select the interaction voltages and duration so that
the resulting transverse evolution is negligible over the pulse.
If this condition cannot be reached, the full coupled Hamiltonian is needed
(Appendix~\ref{sec:calibration-parallel-st-gates}). Under this
condition and with negligible charge leakage, the nonlocal factor
of the interaction pulse is
\begin{equation}
 U_{\mathrm{ent}}(\Phi)=\exp(-i\Phi Z_AZ_B/4),
 \qquad \Phi=\frac{1}{\hbar}\int_0^\tau\kappa(t)\,dt.
 \label{eq:st-entangling-gate}
\end{equation}
To reach $\Phi=\pi$, measure the conditional phase of the full
joint voltage pulse and adjust its barrier and detuning settings
or its duration (Appendix~\ref{sec:calibration-conditional-phase-measurement}).
The local phases accumulated in the same pulse must also be
measured; they determine the net local corrections that convert
the full pulse into CZ. In terms of the isolated nonlocal factor,
the ideal conversion is
\begin{equation}
 \mathrm{CZ}_{AB}
   =e^{-i\pi/4}R_{z,A}(-\pi/2)R_{z,B}(-\pi/2)
       U_{\mathrm{ent}}(\pi),
 \label{eq:st-cz-gate}
\end{equation}
Using the calibrated local rotations to implement Hadamard gates
on $B$ before and after CZ gives the encoded CNOT
[Eq.~\eqref{eq:exchange-cnot-construction}], with $\ket0=\ket S$ and
$\ket1=\ket{T_0}$ in each double dot.
The physical gate acts on four spins.
Uncompensated local $X$ terms do not commute with $Z_AZ_B$ and
require the full time-ordered evolution.
Capacitive singlet--triplet entanglement was demonstrated in
GaAs~\cite{Shulman2012}; silicon requires its own coupling calibration.

\paragraph{Charge admixture sets the conditional energy.}
The conditional energy $\kappa$ comes from the Coulomb interaction
between the two double dots. In a density-interaction model without
inter-block tunneling,
\begin{equation}
 H_{\mathrm{micro}}=H_{\mathrm{Hub},A}+H_{\mathrm{Hub},B}
                  +H_C^{AB},
 \qquad
 H_C^{AB}=\sum_{i,j=1}^{2}W_{ij}n_{A,i}n_{B,j},
 \label{eq:st-interblock-coulomb}
\end{equation}
where $n_{A,i}$ counts electrons on dot $i$ in $A$, and $W_{ij}$
is its Coulomb coupling to dot $j$ in $B$, in energy units.
These cross-qubit terms are additional to each double dot's $U,V$.
Coulomb interaction acts on charge, so it can distinguish the logical
spin states only when they have different charge distributions.
Near a $(1,1)$--$(0,2)$ anticrossing, tunneling mixes
$\ket{S(1,1)}$ with $\ket{S(0,2)}$, while $\ket{T_0}$ remains
predominantly in $(1,1)$. The resulting dipole difference makes the
inter-qubit energy depend on both logical states.

For a singlet with probability $p_q$ of occupying $(0,2)$ and a
triplet confined to $(1,1)$, the weak inter-block density interaction
gives
\begin{equation}
 \kappa=p_Ap_B(W_{11}-W_{12}-W_{21}+W_{22}).
 \label{eq:platform-st-charge-coupling}
\end{equation}
This assumes that the coupling does not appreciably change the
separate charge eigenstates. Gate voltages change $p_A,p_B$ through
detuning and tunneling, thereby controlling the conditional energy.
Increasing charge admixture can strengthen coupling but also
increase charge-noise sensitivity~\cite{Burkard2023,Shulman2012}.
Appendix~\ref{sec:calibration-capacitive-st-interaction} obtains the
admixtures from the charge Hamiltonian and separates the energy
shifts in Eq.~\eqref{eq:st-interaction-decomposition}.

\paragraph{Local gates in the presence of residual coupling.}
The same interaction is present during a nominally local gate if
$\kappa$ remains nonzero. In the diagonal limit above, two local
gates act independently only when the residual $\Phi$ is negligible
or is cancelled, modulo $2\pi$ and local phases. With simultaneous
local $X$ terms, independence must instead be checked from the full
two-qubit evolution. Reducing charge admixture suppresses $\kappa$
but also changes each qubit's exchange; these controls cannot be
assumed to be independent. In the diagonal limit, measure the local
and residual conditional phases of the proposed simultaneous pulses
and adjust their voltages, durations, or scheduling to retain the
target local angles while canceling the unwanted conditional phase
(Appendix~\ref{sec:calibration-parallel-st-gates}). An echoed
entangling pulse is constructed in
Appendix~\ref{sec:calibration-echo-control}.

\subsection{Gates for exchange-only qubits}
\label{sec:exchange-only-gates}
\label{sec:encoded-gates}

\paragraph{Single-pair exchange gates.}
Here the goal is to generate logical rotations using exchange pulses
alone. In the basis of Eq.~\eqref{eq:exchange-only-basis}, exchange
on spins 1 and 2 distinguishes their singlet and triplet, whereas
exchange on spins 2 and 3 mixes the encoded states.
We determine the two rotation axes from the physical Hamiltonian
\begin{equation}
  H_{\mathrm E}=J_{12}\,\mathbf{s}_1\cdot\mathbf{s}_2
                  +J_{23}\,\mathbf{s}_2\cdot\mathbf{s}_3,
  \label{eq:exchange-only-control}
\end{equation}
Isotropic pair exchange conserves total spin and its $z$ component,
so it preserves the $(S,m)=(1/2,+1/2)$ computational space even
when its strength changes in time. Nonuniform magnetic fields
need not preserve that space. In the encoded basis, and after
dropping a common energy, the Hamiltonian is
\begin{equation}
 H_{\mathrm{E,q}}
   =\frac{\sqrt3 J_{23}}{4}X_{\mathrm E}
      +\left(-\frac{J_{12}}2+\frac{J_{23}}4\right)Z_{\mathrm E}.
 \label{eq:exchange-only-logical-control}
\end{equation}
The two projected exchange matrices are evaluated in
Appendix~\ref{sec:calibration-eo-logical-axes},
Eq.~\eqref{eq:exchange-only-projection}.
For a pulse on one pair with the other exchange negligible, let
$\theta_{ij}=\int_0^\tau J_{ij}(t)\,dt/\hbar$.
The active pair's exchange area and the other pair's residual
exchange must be measured at the selected voltage settings
(Appendix~\ref{sec:calibration-eo-pulse-areas}). Choose the compensated
barrier waveform and its duration for the required $\theta_{ij}$,
with the inactive pair held at its calibrated low-exchange setting.
Up to global phase,
\begin{equation}
 \begin{aligned}
 U_{12,\mathrm E}(\theta_{12})
    &=\exp(i\theta_{12}Z_{\mathrm E}/2)
      =R_{z,\mathrm E}(-\theta_{12}),\\
 U_{23,\mathrm E}(\theta_{23})
    &=\exp\!\left[-\frac{i\theta_{23}}2
          \left(\frac{\sqrt3}{2}X_{\mathrm E}
                            +\frac12 Z_{\mathrm E}\right)\right].
 \end{aligned}
 \label{eq:exchange-only-rotations}
\end{equation}
The two pulses rotate about $-\hat z$ and $(\sqrt3/2,0,1/2)$.
Because these axes are nonparallel, their generators and commutators
span the single-qubit rotation algebra.
Choosing pulse areas and alternating the two couplings therefore
gives arbitrary encoded rotations~\cite{DiVincenzo2000}.

\paragraph{A direct logical \texorpdfstring{$R_x$}{Rx} gate.}
For constant simultaneous couplings, $J_{12}=J_{23}/2$ cancels the
$Z_{\mathrm E}$ term. Both couplings must be calibrated with the
two barrier controls applied together; the simultaneous voltages
are adjusted until the measured couplings satisfy this ratio
(Appendix~\ref{sec:calibration-eo-direct-x}). Holding these couplings
for time $\tau$ gives
\begin{equation}
 H_{\mathrm{E,q}}=\frac{\sqrt3 J_{23}}{4}X_{\mathrm E},\qquad
 U_{\mathrm E}(\tau)
    =R_{x,\mathrm E}\!\left(\frac{\sqrt3 J_{23}\tau}{2\hbar}\right).
 \label{eq:exchange-only-direct-x-gate}
\end{equation}
For time-dependent pulses, the same gate follows if the ratio is
maintained at every instant, including the ramps:
\begin{equation}
 J_{12}(t)=\frac12J_{23}(t)
 \quad\Longrightarrow\quad
 U_{\mathrm E}(\tau)=R_{x,\mathrm E}(\vartheta_x),\qquad
 \vartheta_x=\frac{\sqrt3}{2\hbar}\int_0^\tau J_{23}(t)\,dt.
 \label{eq:exchange-only-shaped-x-gate}
\end{equation}
A ratio imposed only during the constant hold does not make the
entire gate a pure $x$ rotation. The complete gate angle must
therefore be calibrated for the chosen pair of waveforms. If the
ratio cannot be maintained on their edges, use a sequence of the
calibrated single-pair rotations instead
(Appendix~\ref{sec:calibration-eo-direct-x}).

\paragraph{Encoded CNOT: target gate and pulse requirements.}
To implement CNOT [Eq.~\eqref{eq:exchange-only-cnot-definition}]
between two triple-dot qubits $A=(1,2,3)$ and $B=(4,5,6)$, we need
exchange between the blocks and a pulse sequence that returns both
blocks to their computational spaces. The logical states in the
target gate are $\ket{0_{\mathrm E}},\ket{1_{\mathrm E}}$ in each block.
For a linear six-dot chain, the available Hamiltonian is
\begin{equation}
 H_6(t)=\sum_{i=1}^{5}J_{i,i+1}(t)
                 \mathbf{s}_i\cdot\mathbf{s}_{i+1}.
 \label{eq:six-spin-exchange}
\end{equation}
The boundary term $J_{34}$ couples the blocks, but it also creates
a leakage problem absent from the one-block gate.
In particular,
$[\mathbf{s}_3\cdot\mathbf{s}_4,\mathbf{s}_A^2]\ne0$, where
$\mathbf{s}_A=\mathbf{s}_1+\mathbf{s}_2+\mathbf{s}_3$.
It can change either block's total spin or its $z$ component despite
conserving total six-spin angular momentum.
Thus a physical $\sqrt{\mathrm{SWAP}}$ on spins 3 and 4 is not an
encoded gate by itself.

The pulse sequence must cancel any final leakage while retaining
an entangling action on the code. To test both requirements,
calculate the product of $M$ pair pulses in the full six-spin space:
\begin{equation}
 U_6=U_{i_Mj_M}(\theta_M)\cdots U_{i_1j_1}(\theta_1),
 \qquad
 U_{ij}(\theta)=\exp[-i\theta\,\mathbf{s}_i\cdot\mathbf{s}_j].
 \label{eq:encoded-six-spin-sequence}
\end{equation}
Let $P_{AB}$ project onto the product of the two code spaces,
each with $S=1/2,m=+1/2$. An encoded CNOT, including local corrections,
must satisfy
\begin{equation}
 \begin{aligned}
  (I-P_{AB})U_6P_{AB}&=0,\\
  \left.P_{AB}U_6P_{AB}\right|_{\mathrm{q},A\otimes\mathrm{q},B}
      &=e^{i\chi}\mathrm{CNOT}_{A\to B},
 \end{aligned}
 \label{eq:encoded-cnot-target}
\end{equation}
Here $\chi$ is a global phase.
The first condition excludes final leakage; the second fixes the
logical gate. Intermediate leakage is allowed, so exponentiating
only the projected boundary exchange is insufficient.

Fong and Wandzura give exact six-spin exchange-only
CNOT sequences~\cite{FongWandzura2011}. Their pulse angles and
ordering are not reproduced here. For a chosen sequence, each
target exchange area must be converted to a barrier waveform and
duration using the measured response of that pair, including the
boundary pair 3--4. The final leakage and logical action of the
complete sequence, including relative phases, must then be measured
(Appendix~\ref{sec:calibration-eo-two-qubit-gates}). Together with
one-block rotations, such a gate gives universal encoded computation.

